\section{Title}

\begin{frame}[plain]
  \vspace*{0.20cm}
  \begin{columns}[c,onlytextwidth]
    \begin{column}{0.58\textwidth}
      {\fontsize{22}{25}\selectfont\bfseries Adversarial Perturbations\\[0.18em] in Social Networks\par}
      \vspace{0.30cm}
      {\fontsize{12}{14}\selectfont\bfseries How small graph edits can change perceived power and influence\par}
      \vspace{0.42cm}
      {\color{PresFourAdversary}\rule{0.10\linewidth}{0.7pt}%
       \color{PresFourFake}\rule{0.10\linewidth}{0.7pt}%
       \color{PresFourDecrease}\rule{0.10\linewidth}{0.7pt}\par}
      \vspace{0.42cm}
      {\small Krzysztof Czerwionka\hspace{0.65cm}Przemysław Fuchs\par}
      \vspace{0.16cm}
      {\small TGS, MIM UW\par}
      \vspace{0.18cm}
      {\footnotesize\color{PresFourMuted}Paper: Avram, Mishra, Parulian, Diesner (ASONAM 2019)\par}
      \vspace{0.04cm}
      {\footnotesize\color{PresFourMuted}Class: game-theoretic approach to social network analysis\par}
    \end{column}
    \begin{column}{0.40\textwidth}
      \centering
      \begin{tikzpicture}
        \node[modelbox,inner sep=8pt] {\TinyGraph[fake][1.05]};
      \end{tikzpicture}
      \vspace{0.20cm}

      {\scriptsize\color{PresFourMuted}red = adversary \quad purple = fake node\\[0.1em]dashed red = added/removed tie}
    \end{column}
  \end{columns}
  \vfill
  \SlideNote{1:00}{Open with the idea that the paper treats social network metrics as something that can be attacked, not just measured.}
\end{frame}
